\documentclass[reqno]{amsart}
\usepackage[T1]{fontenc}
\usepackage{amsmath,amssymb,amsthm}
\usepackage{xr-hyper}
\usepackage[hidelinks]{hyperref}
\externaldocument[gf-]{20261125T005000--galois-fourier-and-averaging__workbook}
\externaldocument[cs-]{20261125T005500--columns-from-a-sphere__workbook}
\externaldocument[sm-]{20261125T000000--slope-method__workbook}
\externaldocument[ch-]{20261125T003000--the-chain-and-the-square-root__workbook}
\externaldocument[fk-]{20261125T009500--fock-model-and-theta-intertwiners__workbook}
\externaldocument[cm-]{20261124T223100--complex-multiplication-and-inert-reduction__workbook}
\externaldocument[od-]{20261125T008000--one-cm-determinant__workbook}
\externaldocument[st-]{20261124T214500--supersingular-theta-divisibility__workbook}
\externaldocument[ip-]{20261124T214600--the-isotropic-plane__workbook}
\externaldocument[hn-]{20261125T013500--hecke-nondegeneracy-on-the-tree__workbook}
\externaldocument[pf-]{20261125T013700--pairings-on-a-formal-group-of-height-two__workbook}
\externaldocument[ho-]{20261125T016500--hecke-operators-and-theta-lifts__workbook}
\externaldocument[dc-]{20261125T016700--divisibility-along-the-chain__workbook}
\externaldocument[nk-]{20261126T020110--the-normalized-kernel-law__workbook}
\externaldocument[fo-]{20261126T020120--frames-at-odd-positions-and-unit-scalars__workbook}
\externaldocument[as-]{20261125T021500--assembly-at-one-scale__workbook}
\externaldocument[ea-]{20261126T020130--nondegeneracy-at-an-edge__workbook}
\externaldocument[co-]{20261126T020140--the-colength-of-a-hecke-image__workbook}
\externaldocument[gl-]{20261126T020150--gain-and-loss-at-an-inert-prime__workbook}

\newcounter{chap}
% Chapter numbers, generated from 20261125T012000--grh-workbook__guide.org
% by scripts/make-chapters.py.  Do not edit: change the order of the
% index and rerun the script (./preview-site.sh and CI do this).
%
% Usage in a sheet, after \newcounter{chap}:
%   % Chapter numbers, generated from 20261125T012000--grh-workbook__guide.org
% by scripts/make-chapters.py.  Do not edit: change the order of the
% index and rerun the script (./preview-site.sh and CI do this).
%
% Usage in a sheet, after \newcounter{chap}:
%   % Chapter numbers, generated from 20261125T012000--grh-workbook__guide.org
% by scripts/make-chapters.py.  Do not edit: change the order of the
% index and rerun the script (./preview-site.sh and CI do this).
%
% Usage in a sheet, after \newcounter{chap}:
%   \input{chapters}\chapnumber{<sheet ID>}
\makeatletter
\providecommand\chapters@set[2]{\@namedef{chapters@#1}{#2}}
\chapters@set{20261125T001000}{1}% zeros-to-ramp-sums
\chapters@set{20261125T001100}{2}% inert-primes-and-estimate-b
\chapters@set{20261125T005000}{3}% galois-fourier-and-averaging
\chapters@set{20261125T005500}{4}% columns-from-a-sphere
\chapters@set{20261125T000000}{5}% slope-method
\chapters@set{20261126T030100}{6}% a-twist-for-each-prime
\chapters@set{20261125T003000}{7}% the-chain-and-the-square-root
\chapters@set{20261125T010000}{8}% rank-lemma-and-zero-estimate
\chapters@set{20261125T010500}{9}% poincare-series-and-near-orthogonality
\chapters@set{20261124T223000}{10}% elliptic-curves-for-the-proof
\chapters@set{20261125T009500}{11}% fock-model-and-theta-intertwiners
\chapters@set{20261124T223100}{12}% complex-multiplication-and-inert-reduction
\chapters@set{20261125T008000}{13}% one-cm-determinant
\chapters@set{20261124T214500}{14}% supersingular-theta-divisibility
\chapters@set{20261124T214600}{15}% the-isotropic-plane
\chapters@set{20261125T013700}{16}% pairings-on-a-formal-group-of-height-two
\chapters@set{20261125T013500}{17}% hecke-nondegeneracy-on-the-tree
\chapters@set{20261126T030200}{18}% the-bound-is-attained
\chapters@set{20261125T009000}{19}% kronecker-limit-and-chowla-selberg
\chapters@set{20261125T002000}{20}% definite-theta-lifts
\chapters@set{20261125T015500}{21}% siegel-weil-in-the-smallest-case
\chapters@set{20261125T002100}{22}% definite-theta-lifts-and-norms
\chapters@set{20261125T016500}{23}% hecke-operators-and-theta-lifts
\chapters@set{20261125T016700}{24}% divisibility-along-the-chain
\chapters@set{20261125T004000}{25}% denominators-and-sha
\chapters@set{20261125T004100}{26}% tate-twists-and-sha
\chapters@set{20261126T020430}{27}% the-minus-part-by-gauss-sums
\chapters@set{20261125T014100}{28}% congruence-modules
\chapters@set{20261125T014500}{29}% eisenstein-congruences-and-selmer-classes
\chapters@set{20261125T014700}{30}% pseudocharacters-and-the-length-bound
\chapters@set{20261125T011000}{31}% galois-representations-black-box
\chapters@set{20261126T030300}{32}% the-data
\chapters@set{20261126T020210}{33}% bessel-poincare-series-and-their-gram-operators
\chapters@set{20261126T030400}{34}% the-seesaw-identity
\chapters@set{20261126T030500}{35}% divisibility-of-the-composite-lift
\chapters@set{20261126T030600}{36}% theorem-a-short-forms
\chapters@set{20261126T030700}{37}% the-determinant-of-a-chain
\chapters@set{20261126T030800}{38}% the-cost-of-one-lift
\chapters@set{20261126T030900}{39}% denominators-of-projections-to-lifts
\chapters@set{20261126T031000}{40}% twists-with-few-congruences
\chapters@set{20261126T031100}{41}% theorem-c-and-the-main-inequality
\providecommand\chapnumber[1]{%
  \@ifundefined{chapters@#1}%
    {\PackageError{chapters}{Sheet #1 is not listed in the index}%
      {List it in 20261125T012000--grh-workbook__guide.org\MessageBreak
       and run scripts/make-chapters.py.}}%
    {\edef\thechap{\@nameuse{chapters@#1}}}}
\makeatother
\chapnumber{<sheet ID>}
\makeatletter
\providecommand\chapters@set[2]{\@namedef{chapters@#1}{#2}}
\chapters@set{20261125T001000}{1}% zeros-to-ramp-sums
\chapters@set{20261125T001100}{2}% inert-primes-and-estimate-b
\chapters@set{20261125T005000}{3}% galois-fourier-and-averaging
\chapters@set{20261125T005500}{4}% columns-from-a-sphere
\chapters@set{20261125T000000}{5}% slope-method
\chapters@set{20261126T030100}{6}% a-twist-for-each-prime
\chapters@set{20261125T003000}{7}% the-chain-and-the-square-root
\chapters@set{20261125T010000}{8}% rank-lemma-and-zero-estimate
\chapters@set{20261125T010500}{9}% poincare-series-and-near-orthogonality
\chapters@set{20261124T223000}{10}% elliptic-curves-for-the-proof
\chapters@set{20261125T009500}{11}% fock-model-and-theta-intertwiners
\chapters@set{20261124T223100}{12}% complex-multiplication-and-inert-reduction
\chapters@set{20261125T008000}{13}% one-cm-determinant
\chapters@set{20261124T214500}{14}% supersingular-theta-divisibility
\chapters@set{20261124T214600}{15}% the-isotropic-plane
\chapters@set{20261125T013700}{16}% pairings-on-a-formal-group-of-height-two
\chapters@set{20261125T013500}{17}% hecke-nondegeneracy-on-the-tree
\chapters@set{20261126T030200}{18}% the-bound-is-attained
\chapters@set{20261125T009000}{19}% kronecker-limit-and-chowla-selberg
\chapters@set{20261125T002000}{20}% definite-theta-lifts
\chapters@set{20261125T015500}{21}% siegel-weil-in-the-smallest-case
\chapters@set{20261125T002100}{22}% definite-theta-lifts-and-norms
\chapters@set{20261125T016500}{23}% hecke-operators-and-theta-lifts
\chapters@set{20261125T016700}{24}% divisibility-along-the-chain
\chapters@set{20261125T004000}{25}% denominators-and-sha
\chapters@set{20261125T004100}{26}% tate-twists-and-sha
\chapters@set{20261126T020430}{27}% the-minus-part-by-gauss-sums
\chapters@set{20261125T014100}{28}% congruence-modules
\chapters@set{20261125T014500}{29}% eisenstein-congruences-and-selmer-classes
\chapters@set{20261125T014700}{30}% pseudocharacters-and-the-length-bound
\chapters@set{20261125T011000}{31}% galois-representations-black-box
\chapters@set{20261126T030300}{32}% the-data
\chapters@set{20261126T020210}{33}% bessel-poincare-series-and-their-gram-operators
\chapters@set{20261126T030400}{34}% the-seesaw-identity
\chapters@set{20261126T030500}{35}% divisibility-of-the-composite-lift
\chapters@set{20261126T030600}{36}% theorem-a-short-forms
\chapters@set{20261126T030700}{37}% the-determinant-of-a-chain
\chapters@set{20261126T030800}{38}% the-cost-of-one-lift
\chapters@set{20261126T030900}{39}% denominators-of-projections-to-lifts
\chapters@set{20261126T031000}{40}% twists-with-few-congruences
\chapters@set{20261126T031100}{41}% theorem-c-and-the-main-inequality
\providecommand\chapnumber[1]{%
  \@ifundefined{chapters@#1}%
    {\PackageError{chapters}{Sheet #1 is not listed in the index}%
      {List it in 20261125T012000--grh-workbook__guide.org\MessageBreak
       and run scripts/make-chapters.py.}}%
    {\edef\thechap{\@nameuse{chapters@#1}}}}
\makeatother
\chapnumber{<sheet ID>}
\makeatletter
\providecommand\chapters@set[2]{\@namedef{chapters@#1}{#2}}
\chapters@set{20261125T001000}{1}% zeros-to-ramp-sums
\chapters@set{20261125T001100}{2}% inert-primes-and-estimate-b
\chapters@set{20261125T005000}{3}% galois-fourier-and-averaging
\chapters@set{20261125T005500}{4}% columns-from-a-sphere
\chapters@set{20261125T000000}{5}% slope-method
\chapters@set{20261126T030100}{6}% a-twist-for-each-prime
\chapters@set{20261125T003000}{7}% the-chain-and-the-square-root
\chapters@set{20261125T010000}{8}% rank-lemma-and-zero-estimate
\chapters@set{20261125T010500}{9}% poincare-series-and-near-orthogonality
\chapters@set{20261124T223000}{10}% elliptic-curves-for-the-proof
\chapters@set{20261125T009500}{11}% fock-model-and-theta-intertwiners
\chapters@set{20261124T223100}{12}% complex-multiplication-and-inert-reduction
\chapters@set{20261125T008000}{13}% one-cm-determinant
\chapters@set{20261124T214500}{14}% supersingular-theta-divisibility
\chapters@set{20261124T214600}{15}% the-isotropic-plane
\chapters@set{20261125T013700}{16}% pairings-on-a-formal-group-of-height-two
\chapters@set{20261125T013500}{17}% hecke-nondegeneracy-on-the-tree
\chapters@set{20261126T030200}{18}% the-bound-is-attained
\chapters@set{20261125T009000}{19}% kronecker-limit-and-chowla-selberg
\chapters@set{20261125T002000}{20}% definite-theta-lifts
\chapters@set{20261125T015500}{21}% siegel-weil-in-the-smallest-case
\chapters@set{20261125T002100}{22}% definite-theta-lifts-and-norms
\chapters@set{20261125T016500}{23}% hecke-operators-and-theta-lifts
\chapters@set{20261125T016700}{24}% divisibility-along-the-chain
\chapters@set{20261125T004000}{25}% denominators-and-sha
\chapters@set{20261125T004100}{26}% tate-twists-and-sha
\chapters@set{20261126T020430}{27}% the-minus-part-by-gauss-sums
\chapters@set{20261125T014100}{28}% congruence-modules
\chapters@set{20261125T014500}{29}% eisenstein-congruences-and-selmer-classes
\chapters@set{20261125T014700}{30}% pseudocharacters-and-the-length-bound
\chapters@set{20261125T011000}{31}% galois-representations-black-box
\chapters@set{20261126T030300}{32}% the-data
\chapters@set{20261126T020210}{33}% bessel-poincare-series-and-their-gram-operators
\chapters@set{20261126T030400}{34}% the-seesaw-identity
\chapters@set{20261126T030500}{35}% divisibility-of-the-composite-lift
\chapters@set{20261126T030600}{36}% theorem-a-short-forms
\chapters@set{20261126T030700}{37}% the-determinant-of-a-chain
\chapters@set{20261126T030800}{38}% the-cost-of-one-lift
\chapters@set{20261126T030900}{39}% denominators-of-projections-to-lifts
\chapters@set{20261126T031000}{40}% twists-with-few-congruences
\chapters@set{20261126T031100}{41}% theorem-c-and-the-main-inequality
\providecommand\chapnumber[1]{%
  \@ifundefined{chapters@#1}%
    {\PackageError{chapters}{Sheet #1 is not listed in the index}%
      {List it in 20261125T012000--grh-workbook__guide.org\MessageBreak
       and run scripts/make-chapters.py.}}%
    {\edef\thechap{\@nameuse{chapters@#1}}}}
\makeatother

% Use the table as soon as the index includes this sheet.
\expandafter\providecommand\csname chapters@20261125T020100\endcsname{?}
\chapnumber{20261125T020100}
\renewcommand{\thesection}{\thechap.\arabic{section}}
\theoremstyle{definition}
\newtheorem{problem}{Problem}[section]
\theoremstyle{plain}
\newtheorem{theorem}{Theorem}[section]
\newcommand{\Z}{\mathbb Z}
\newcommand{\Q}{\mathbb Q}
\newcommand{\Oc}{\mathcal O}
\newcommand{\Int}{\mathrm{Int}}
\newcommand{\capn}{\mathrm{cap}}
\newcommand{\Sig}{\mathrm{Sig}}
\newcommand{\io}{\mathrm{io}}
\newcommand{\IO}{\mathrm{IO}}
\newcommand{\LP}{\mathrm{LP}}
\setlength{\emergencystretch}{2em}
\title{Divisibility at inert primes, in general}
\author{}
\date{November 25, 2026}
\begin{document}
\maketitle
\refstepcounter{chap}\label{chap}

\paragraph{Prerequisites.}
Chapters~\ref{cm-chap}, \ref{st-chap} and \ref{ip-chap} on inert CM reduction and invariant derivatives;
Chapter~\ref{pf-chap} on formal pairings; Chapters~\ref{hn-chap} and \ref{dc-chap} on the integral Hecke image and the composite of lifts.  The simplest incarnations below recall Chapter~\ref{gf-chap} (the prime sum), Chapter~\ref{cs-chap} (the tops), Chapters~\ref{sm-chap} and~\ref{ch-chap} (the slope method and the chain), Chapter~\ref{fk-chap} (theta intertwiners), Chapter~\ref{od-chap} (one CM determinant) and Chapter~\ref{ho-chap} (Hecke operators and theta lifts).  We use the convention \(v_p(p)=1\), \((z)_+=\max(z,0)\).  This chapter states results of \cite[\S\S8--12]{Adr} in the form used later and checks them on examples; it does not reproduce their definitions in full, and the reader needs \cite[\S\S8--12]{Adr} open.

Which powers of \(p\) gained by differentiating a theta kernel survive the chain of lifts, and what is the index of the source on which they survive?

\section{The scalar in the kernel law}\label{sec:kernel}
In the manuscript \(K=\Q(\sqrt{-d})\), where the auxiliary prime \(d\equiv7\bmod16\); the CM period lattice \(\Lambda\) has discriminant one. Fix \(p\ge5\) inert in \(K\) and the good integral model of \cite[\S8.2]{Adr}. Write \(W=W(\overline{\mathbb F}_p)\), with Frobenius \(\sigma\).

\begin{theorem}[Lubin--Tate coordinate]\label{thm:LT}
Over \(W\) the curve's formal \(\Oc_{K,p}\)-module admits an integral coordinate with unit linear term and logarithm \(\operatorname{lg}(X)=\sum_{k\ge0}X^{p^{2k}}/p^k\) \cite{LT}, \cite[Prop.~8.11]{Adr}.
\end{theorem}
\noindent \emph{Status}: classical. Simplest incarnation: Problem~\ref{cm-pr:model-log} gives an integral change of parameter for the Gaussian curve, with uniformizer \(-p\).

\begin{theorem}[Kernel law]\label{thm:kernel}
Let \(M\) be a positive definite hermitian \(\Oc_K\)-lattice, \(L_M=M\otimes_{\Oc_K}\Lambda\), and \(NL_M\subset L\subset N^{-1}L_M\) with \(p\nmid N\). Let \(G\) be a theta function on \(L\) for the induced polarization, with sign semicharacter and algebraic Taylor coefficients. In an orthonormal frame of \(M_p\), the Taylor series corrected by the product of the one-variable exponential factors of \cite[Thm.~8.19]{Adr} has bounded coefficients. Suppose \(M_p\) is an orthogonal sum of hyperbolic planes and a lattice with orthonormal basis. Write \(x_e,y_e\) for the Lubin--Tate coordinates of the two isotropic vectors in each plane. The restriction to these planes in the period normalization of \cite[Thm.~8.21]{Adr} has the form
\[
 \Theta(x,y)=c\exp\!\left(-\sum_e\beta(x_e,y_e)\right)U(x,y),
 \qquad c\in\mathbb C_p^\times,\quad U\in\Oc_{\mathbb C_p}[[x,y]] .
\]
Here \(\beta\) is the coboundary of the one-variable theta correction in \cite[Eq.~(8.6)]{Adr}.
\end{theorem}
\noindent \emph{Status}: proved in \cite{Adr}.  The theorem says nothing about the scalar \(c\).  For the kernels of the Lean formalization \cite{Lean}, built from one theta function of the CM curve normalized to take the value \(1\) at the origin, it holds with \(c=1\) (Theorem~\ref{nk-thm:B} of Chapter~\ref{nk-chap}).  Simplest incarnation: Theorem~\ref{ip-thm:line} gives an integral restriction \(\widehat P(X)\) to one isotropic line for \(K=\Q(i)\), \(p\equiv3\bmod8\). Theorem~\ref{st-thm:single} gives \(v_p(D^n f(0))\ge\sum_{i>0,\ i\ {\rm odd}}\lfloor n/p^i\rfloor\) for an integral series \(f\) and its invariant derivation \(D\).

\begin{problem}\label{pr:line}
In one hyperbolic plane \(\beta(x,0)=0\). Show that \(c^{-1}\Theta(x,0)\) is integral. In the Gaussian model with its integral line series \(\widehat P\), show that \(v_p(D^{p^3}\widehat P(0))\ge p^2+1\) for \(p\equiv3\bmod8\).
\end{problem}

\section{What a minor can absorb}\label{sec:minors}
Index the planes by \(0\le r<\nu,\ 0\le a<r_0,\ 0\le b<r_1\). The \emph{positions} \(s\in\Z/(2\nu)\) number the places of the CM type above \(p\) in one Frobenius orbit, so that Frobenius followed by complex conjugation moves the place at \(s\) to the place at \(s+1\) \cite[\S11.1]{Adr}; Section~\ref{fo-sec:positions} of Chapter~\ref{fo-chap} gives the details.  After the trace-dual change of coordinates, the variable matrices \(X_s=(X_{s,a,b})\) are indexed by the positions. For \(\varpi_p^{p^2-1}=p\), put
\[
 \mathsf S_{s,a,b}=\sum_{k\ge0}(-1)^k X_{s-2k,a,b}^{p^{2k}},
 \qquad h_\upsilon=\sum_{s,a,b}\upsilon_s\mathsf S_{s,a,b}^p\mathsf S_{s+1,a,b}.
\]
For \(h(0)=0\), let \(\mathcal C(h)\) consist of the series \(\sum_{i\ge0}p^{-i}h^ig_i\) with integral \(g_i\). Each denominator \(p^i\) therefore comes with \(i\) factors of \(h\).

\begin{theorem}[Trace-dual coordinates]\label{thm:class}
In Theorem~\ref{thm:kernel}, use trace-dual bases \(\omega_r,\omega_r^*\) of the unramified degree-\(\nu\) extension of \(\Q_{p^2}\) over \(\Q_{p^2}\), integral on both sides. In the two logarithmic variables substitute respectively
\(\sum_{s\ {\rm even}}\sigma^s(\omega_r^*)X_{s,a,b}\) and
\(\sum_{s\ {\rm odd}}\sigma^s(\omega_r)X_{s,a,b}\); call the resulting series \(\Theta^\flat\).
Then \(c^{-1}\Theta^\flat(\varpi_p X)\in\mathcal C(h_\upsilon)\) for some integral \(\upsilon_s\). Indeed, write \(A,B\) for those two linear forms, and \(\phi\) for the compositional inverse of \(\varpi_p^{-1}\operatorname{lg}(\varpi_p X)\). For \(\sigma^2 z=z\) in \(W\), the congruence of \cite[Thm.~9.4]{Adr} is
\[
 \sum_{r,a,b}\big(z\phi(A)\phi(B)^p-\sigma(z)\phi(B)\phi(A)^p\big)
 \equiv h_\upsilon\pmod p ,
\]
where \(\upsilon_s=z\) for odd \(s\) and \(-\sigma(z)\) for even \(s\). Applying this to the kernel pairing gives the class assertion \cite[Prop.~9.5, Cor.~8.22]{Adr}.
\end{theorem}
\noindent \emph{Status}: proved in \cite{Adr}. Simplest incarnation: Theorem~\ref{pf-thm:class} puts the rescaled one-pair exponential times an integral series in \(\mathcal C(\sigma(c_0)x^py-c_0xy^p)\), where \(c_0\) is its pairing parameter.

\begin{problem}\label{pr:odd}
Work modulo \(p\). In \(\mathsf S_{s,a,b}^p\mathsf S_{s+1,a,b}\), take the \(k\)-th term from the first factor after Frobenius and the linear term from the second. Its exponents are \(p^{2k+1}\) and \(1\); its positions differ by \(2k+1\) modulo \(2\nu\).
\end{problem}

A \emph{shape} consists of a set \(\mathcal T\) of positions, called \emph{tops}, and at each position a list of columns of heights \(h_c\le r_0\), with a list \(I_c\) of \(h_c\) column indices for each \cite[\S9]{Adr}. Use the first \(h_c\) rows of a common frame \(Y_sX_s\) and the columns \(I_c\) to form a determinant; \(P_I\) is their product over columns and positions. Pair by \(\langle P_I,F\rangle=P_I(\partial)F(0)\), a polynomial in the frame entries.

\begin{problem}\label{pr:cap}
For a square minor of \(X+tuv^{\mathsf T}\), the coefficient of \(t^d\) vanishes for \(d\ge2\). Deduce that a product of \(w\) framed minors has degree at most \(w\) along any rank-one direction, whatever their heights.
\end{problem}

Problem~\ref{pr:cap} explains why each remaining column costs one. At a top \(s\), count its columns of height one as \(\alpha_s\); set \(\alpha_s=0\) at the other positions.  The \emph{remaining} columns at \(s\) are the others: all columns if \(s\) is not a top, and the columns of height at least two if it is.  Write \(\capn_s\) for their number and \(\lambda_s\) for their number of boxes, the sum of their heights; thus \(\lambda_s\) does not count the columns of height one at a top. Put \(\operatorname{len}_p(0)=0\), and let \(\operatorname{len}_p(n)\) count base-\(p\) digits for \(n>0\). For a top \(s\), let \(E(s)\) consist of the least nonnegative residues of \(s'-s\) for tops \(s'\) of opposite parity. For positive odd \(i\) put \(A_{s,i}=(\alpha_s/p^i-\capn_{s+i})_+\). The bounds of \cite[Thm.~9.6]{Adr} use
\begin{align*}
 \Delta&=\sum_s\big(r_1\operatorname{len}_p(\alpha_s)+r_0r_1\operatorname{len}_p(\lambda_s)\big)+\sum_{s\in\mathcal T}\operatorname{len}_p(\alpha_s)^2,\\
 g^A&=\sum_{s\in\mathcal T}\sum_{i\notin E(s)} A_{s,i}-\Delta,\\
 g^B&=\frac p{p^2-1}\sum_{s\in\mathcal T_{\rm ev}}\alpha_s
       +\sum_{s\in\mathcal T_{\rm odd}}\sum_i A_{s,i}-\Delta .
\end{align*}

\begin{theorem}[Valuations of entries]\label{thm:valuation}
If \(F(\varpi_p X)\in\mathcal C(h_\upsilon)\) with integral \(\upsilon_s\), every frame coefficient \(z\) of \(\langle P_I,F\rangle\) satisfies
\[
 v_p(z)\ge g^A,\qquad v_p(z)\ge g^B-\sum_{s\in\mathcal T_{\rm ev}}\alpha_s,
 \qquad v_p(p^{N_J}z)\ge g^B
\]
for every set \(J\) of column indices. Here \(N_J\) counts height-one top columns labelled outside \(J\) at even tops or inside \(J\) at odd tops, and adds, for each remaining column at an odd position, the number of its distinct labels in \(J\) \cite[Thm.~9.6]{Adr}.
\end{theorem}
\noindent \emph{Status}: proved in \cite{Adr}. Simplest incarnation: Theorem~\ref{ip-thm:main} gives the directional bound \(\sum_i(\lfloor a/p^i\rfloor-w)_+\) on the Gaussian plane; Theorem~\ref{pf-thm:two} gives an entry bound both before and after multiplication by \(p^w\) for one formal pair.

\begin{problem}\label{pr:digits}
Use the Gaussian plane for comparison and take \(p=11,\ \nu=r_0=r_1=1\), one even top with \(a=2\cdot11^3\) columns and one column at the other position. The directional bound of Theorem~\ref{ip-thm:main} is \(242\). Show that the two contributing levels in \(g^A+\Delta\) give the same \(242\), while the first inequality of Theorem~\ref{thm:valuation} gives \(221\), since \(\lambda_s\) at the top is zero.
\end{problem}

\section{What survives composition}\label{sec:chain}
A chain datum \cite[Def.~10.1]{Adr} gives integral lattices of equivariant functions \(\Int_j\) on lattice classes; its steps sum local kernel maps over classes with reciprocal stabilizer weights. Let \(d_j\) bound the denominator valuation of those weights at level \(j\).

\begin{theorem}[Integral frames and stabilizers]\label{thm:frames}
For every commutative ring the canonical bideterminants of a fixed shape span a free module with standard-tableau basis, stable under integral row changes and compatible with base change (Green; De Concini--Eisenbud--Procesi). A finite subgroup of \(\mathrm{GL}_n(\Z)\) has order dividing Minkowski's \(M(n)\), where \(v_p M(n)=\sum_{h\ge0}\lfloor n/(p^h(p-1))\rfloor\) at odd \(p\) \cite[Thms.~11.2 and~11.13]{Adr}.
\end{theorem}
\noindent \emph{Status}: classical. Simplest incarnation: the monomials \(x^{k-j}y^j\) give the integral polynomial lattice in Theorem~\ref{hn-thm:vertex}; the reciprocal stabilizer orders occur already in the theta sum of Chapter~\ref{ho-chap}, after Theorem~\ref{ho-thm:eichler}.

\begin{problem}\label{pr:minkowski}
For \([E:\Q]=2f\) and ranks \(N_j\le N_m\), the common denominator \(D_\Gamma=M(2fN_m)\) satisfies \(v_p(D_\Gamma)\le 2fN_m p/(p-1)^2\). For \(p=5,\ 2fN_m=20\) its valuation is \(6\).
\end{problem}

Write \(\mathcal V\) for the finite set of Frobenius orbits. At orbit \(v\) and step \(j\), \(g^0,g^1\) are the two entry bounds and \(\mathrm{Pr}\) is the price of commuting the normalized Hecke operator \(y_v\) through the composite. More precisely, the nontrivial Hecke terms have bound \(g^1\), the trivial term has bound \(\max(g^0,g^1-\alpha)\); the commutation polynomial has integral coefficients and constant coefficient of valuation at least \(\alpha\). The price bounds the valuation of the scalar in this commutation \cite[Defs.~10.4 and~10.5, Thm.~10.7]{Adr}.

\begin{theorem}[Gain on the source]\label{thm:chain}
For a chain datum with the Hecke commutation data and orbit bounds of \cite[Thm.~10.7]{Adr}, choose a set \(B\subseteq\mathcal V\) of orbits, the \emph{served} orbits, and numbers \(S_v\ge\sum_{j=1}^m(g^1-\mathrm{Pr}-g^0)_+(v,j)\) on \(B\). The composite maps the span of
\[
 \left(\prod_{v\in B'}y_v\right)^m p^{\sum_{v\in B\setminus B'} S_v}\Int_0
 \quad(B'\subseteq B)
\]
into \(p^{g^{\rm tot}}\Int_m\), where \(p^z\) on a lattice denotes the valuation filtration at \(z\), and
\[
 g^{\rm tot}=\sum_{j=1}^m\left(\sum_{v\in B}\max(g^0,g^1-\mathrm{Pr})(v,j)
                  +\sum_{v\notin B}g^0(v,j)\right)-\sum_{j=0}^{m-1}d_j .
\]
\end{theorem}
\noindent \emph{Status}: proved in \cite{Adr}.  In the formalization the gain on the Hecke image is Theorem~\ref{gl-thm:gain} of Chapter~\ref{gl-chap}, proved from an input on the lifts that their construction supplies.  Simplest incarnation: Theorem~\ref{dc-thm:chain-gain} gives the sum of the stepwise maxima \(\max(\gamma_j(0),A+\delta_j)\) on \(y^m\Int_0\) for one orbit.

\begin{problem}\label{pr:price}
Serve one orbit for two steps. Take \(d_j=g^0=0\), \(g^1=(3,2)\), \(\mathrm{Pr}=(1,3)\). Show that the composite maps both \(y^2\Int_0\) and \(p^2\Int_0\) into \(p^2\Int_2\), with the whole improvement coming from the first step.
\end{problem}

\begin{theorem}[True lifts]\label{thm:true-lifts}
Let \(p\ge5\) be inert in \(K\) and unramified in the CM field \(E\) of the lift data of \cite[Defs.~11.5, 11.7 and~11.10]{Adr}.  Suppose that
\begin{enumerate}
\item[(a)] every step kernel satisfies Theorem~\ref{thm:kernel} with \(c=1\): in the hyperbolic coordinates there, it is \(\exp(-\sum_e\beta(x_e,y_e))\,U\) with \(U\) integral;
\item[(b)] the scalar normalizing each step changes only by \(p\)-adic units when the lattice at one orbit changes;
\item[(c)] at the odd positions the frames are the duals of the frames at the even positions, and entry estimates pass to inclusions of lattices there as at the even positions, with determinant factors of the opposite sign.
\end{enumerate}
Then the true lifts form a chain datum on the integral frame lattices with \(d_j=v_p(D_\Gamma)\), commuting normalized Hecke operators, and orbit bounds \((g^A,g^B,\alpha_{\rm ev})\), where \(\alpha_{\rm ev}\) counts the height-one columns at even tops. Theorem~\ref{thm:chain} applies with the prices of \cite[Prop.~11.32, Thm.~11.45]{Adr}.
\end{theorem}
\noindent \emph{Status}: derived in \cite[Thm.~11.45]{Adr} from three hypotheses on the lifts, which are (a)--(c) in the form given there: (N3), \cite[Hyp.~11.42]{Adr}, for (a); (N1), \cite[Hyp.~11.22]{Adr}, for (b); and (N2), \cite[Hyp.~11.38]{Adr}, for (c).  Of these, (N1) is proved in \cite[Prop.~11.27]{Adr} for the step normalization of \cite[Def.~11.23]{Adr}, and (N2) and (N3) are assumed there.  The Lean formalization \cite{Lean} proves the final theorem through statements of the form (a)--(c), which Chapters~\ref{nk-chap} and~\ref{fo-chap} develop.  Its kernels are built from one theta function of the CM curve, normalized to take the value \(1\) at the origin, and for them (a) is Theorem~\ref{nk-thm:B}.  Its scalars are \(p\)-adic units, which gives (b) (Section~\ref{fo-sec:N1}).  Its frames at the odd positions are the duals of those at the even positions, and Section~\ref{fo-sec:N2} carries out for them the bookkeeping of (c) that \cite[Rem.~11.39]{Adr} leaves to be written.  We have not checked that the chain datum of the formalization is the one of \cite[Thm.~11.45]{Adr}.  Simplest incarnation: Theorem~\ref{ho-thm:eichler} expresses a sum over neighbours as a Hecke translate plus a scalar multiple of the original theta series; Chapter~\ref{fk-chap} composes the two theta intertwiners in Theorem~\ref{fk-thm:compose}.

Theorem~\ref{thm:kernel} alone does not give (a), since it allows an unknown scalar for each pair of lattice tuples.
\begin{problem}\label{pr:normalization}
Multiplying a kernel with corrected constant term \(1\) by \(p^{-1}\) preserves boundedness in Theorem~\ref{thm:kernel} and prevents any expression with both \(c\) and \(U\) integral. Deduce that applying the boundedness theorem separately to every tuple does not give (a) of Theorem~\ref{thm:true-lifts}.
\end{problem}

\section{The price of passing to the Hecke image}\label{sec:image}
From here on \(x\) is the scale, \(0<\varepsilon<1/52\), \(y=x^\varepsilon\), \(b=e^y\), \(a=xb\), \(k\) is the rank and \(n\) the number of columns, as in \cite[\S5]{Adr}.  Put \(X'=\exp(x^{\varepsilon/3})\), \(b'=a/X'\), \(X''=xnX'\) and \(p_*=k^2x^{26}\) \cite[\S12.1]{Adr}.

For a representation \(V\) of highest weight data \(\lambda\) at a single orbit and \(c_*\ge4(D+1)\), \(\LP(r,\lambda;t,D,\kappa,c_*)\) asks for the following at every target lattice. Choose incoming \(r\)-edges and integral source vectors in their prescribed depth-\(t\) lattices. Each vector vanishes to order \(c_*\) in the edge lattice of every other edge out of its source (order measured in the local uniformizer). Integral combinations of their pairings with tests in the depth-\(t\) target lattice must give \(p^D I_n\) on \(n\ge(1-\kappa)\dim V\) independent tests \cite[Def.~12.6]{Adr}. Thus \(\kappa\) measures a fraction of the dimension.

\begin{theorem}[Hecke nondegeneracy]\label{thm:LP}
Let \(p\) be a prime of \emph{class~I}, meaning that odd Frobenius powers from a top land in the top half of the places \cite[\S12.1]{Adr}, whose Frobenius orbits on the columns have length \(n_p\ge x^{\varepsilon/4}\).  Let \(t_0\) be such an orbit that is \emph{good}, that is, satisfies the spacing conditions of \cite[Def.~5.19]{Adr} on the tops along it, and let \(N(t_0)\) be the number of tops on it.  Assume the size conditions of \cite[\S12.6]{Adr} and the three engines of \cite[Hyp.~12.12]{Adr}: (A) sufficient rank of the reduced segment symbols, (B) detection by the twisted Fock pairing at the stated depth, (C) the mean bound for bad segment shares.  Then for large \(x\) the following holds at every target lattice, with \(r\), \(t\), \(D\) and the share \(\kappa=\kappa(a,b)\) of \cite[Thm.~12.11]{Adr}, where \((a,b)\) are the weights.
\begin{enumerate}
\item[(i)] If \(13\le p<p_*\), then \(\LP(r,\lambda;\lceil t\rceil,D,\kappa,8(D+1))\) holds with \(\kappa\le N(t_0)x^{-12}\) for all weights in the windows.
\item[(ii)] If \(p_*\le p<X'\) and no power \(p^i\) lies in the band \(\lvert p^i-X'\rvert\le x^{-2}X'\), let \(\kappa\) be the least share for which \(\LP(1,\lambda;\lceil t\rceil,D,\kappa,8(D+1))\) holds.  Its mean over the weights in the windows is at most \(N(t_0)x^{-12}\).
\end{enumerate}
\end{theorem}
\noindent \emph{Status}: derived in \cite[Thm.~12.11, Cor.~12.13]{Adr} from the engines \cite[Hyp.~12.12]{Adr}, assumed there.  The proof in \cite[Thm.~12.11]{Adr} is an outline; it also calls on a Pieri splitting and a passage from segment shares to \(\kappa\).  The formalization proves the final theorem; the three engines are not separate statements in it.  Chapter~\ref{ea-chap} states its condition at one edge, which plays the role of \(\LP\), and the theorems that establish it for small primes and for primes of class~I (Theorems~\ref{ea-thm:small} and~\ref{ea-thm:classI}), whose proofs are large and are not developed there; we have not matched the engines to its lemmas.  Simplest incarnation: at one tree vertex, Theorem~\ref{hn-thm:vertex} diagonalizes the image on degree-\(k\) polynomials. For residue size \(Q>k\) the exponent of \(x^{k-j}y^j\) is
\[
 \varepsilon_j=\min\left(k-j,\ \min_{0\le i\le j}
              \left(i+v_p{\binom{k-i}{j-i}}\right)\right).
\]
\begin{problem}\label{pr:vertex}
For \(p=k=5,\ Q=25\), the six exponents are \(0,1,1,1,1,0\). Deduce that the image has rank two after reduction in the target modulo \(5\), but its integral colength over the unramified valuation ring is only four. For every \(Q>k\) this colength is at most \(\sum_{j=0}^k v_p\binom kj\).
\end{problem}

For each candidate prime, first sum \(\Sig(p)=m\sum_{v\ {\rm good}}\kappa_v\) over all its good orbits. The candidates are the class I primes \(p<X'\) with \(n_p\ge x^{\varepsilon/4}\), subject to the power-band exclusion of Theorem~\ref{thm:LP}(ii) for \(p\ge p_*\).  Here \(\kappa_v\) is the share of Theorem~\ref{thm:LP} at the orbit \(v\), and \(t_v\), \(D_p\) are its exponents.  Serve the good orbits only when \(\Sig(p)\le x^{-\varepsilon}\); otherwise the served set \(\mathcal V_p\) is empty. Noncandidates have empty served set as well.

Write \(P(M)\) for the sum of the Hecke images with scalar alternatives in Theorem~\ref{thm:chain}, and \(M^{\rm ram}=M\cap V^{\rm ram}\) for the saturated lattice in the chosen ramified subspace. For \(\mathcal V_p\ne\varnothing\), put
\[
 \tau(p)=e_p+m\sum_{v\in\mathcal V_p}(\lceil t_v\rceil+D_p+1),\quad
 \io(p)=2\Sig(p)\max_{v\in\mathcal V_p} S_v .
\]
Here \(e_p\) is the exponent comparing the two integral structures. For an empty served set put \(\tau(p)=\io(p)=0\); a rejected candidate retains its original \(\Sig(p)\).

\begin{theorem}[Colength]\label{thm:colength}
At \(w\mid p\) for a served class I prime \(p\le X''\), \(p\ne T_u\), take the lattices and Hecke operators at \(w\) given by the chain datum at level zero \cite[\S12.3, (F1)--(F5)]{Adr}, and assume \(\LP\) at every served orbit with parameters \(\lceil t_v\rceil,D_p,\kappa_v,8(D_p+1)\), and \(\Sig(p)\le1/2\).
Then \(N\!P_w=P(M)\cap p^{\lceil\tau(p)\rceil}M\cap V^{\rm ram}\) lies in \(p^{\lceil\tau(p)\rceil}M^{\rm ram}\), and
\[
 \operatorname{length}_{\Oc_w}(M^{\rm ram}/N\!P_w)
 \le v_w(p)\operatorname{rk}(M)\big(\tau(p)+\io(p)\big)
\]
with \(v_w\) the uniformizer valuation \cite[Thm.~12.18]{Adr}. For an empty served set take \(N\!P_w=M^{\rm ram}\).
\end{theorem}
\noindent \emph{Status}: proved in \cite{Adr}.  The formalization proves the colength bound in the form of Theorem~\ref{co-thm:colength} of Chapter~\ref{co-chap}, from the condition at one edge.  Simplest incarnation: Theorem~\ref{hn-thm:vertex} charges the sum of the exponents on a single vertex. Theorem~\ref{sm-thm:one} charges the logarithmic index of a source sublattice divided by its rank.
\begin{problem}\label{pr:share}
Use the inclusion and length bound as a numerical model on a lattice of rank eight, with full ramified part, \(v_w(p)=2\), \(m=2\), one served orbit and \(S_v=3\). For \(\kappa_v=1/8\) and integral \(\tau\), show that the length lies between \(16\tau\) and \(16\tau+24\). With \(\kappa_v=0\), deduce that \(N\!P_w=p^\tau M^{\rm ram}\).
\end{problem}

\section{The prime sum and the slope}\label{sec:global}
Let \(T\) be the set of \(N_T\) tops (Section~\ref{cs-sec:tops}), \(m\) the chain length, \(a_t\) the top weights, \(w_j\) the bottom weights. In the notation of Section~\ref{gf-sec:balance},
\[
 S=\sum_{t\in T}\sum_{j=1}^m P_{t,w_j}(a_t),\qquad
 P_{t,w}(a)=\sum_r\Lambda_d(r)\big(a/r-w(\operatorname{cell}_t(r))\big)_+ .
\]
In particular an admitted unramified inert prime contributes \(2\log p\) times the sum of these positive parts at its odd powers. Let \(\mathrm{gain}^{\mathcal D}_j(p)\) be the sum of the local maxima of Theorem~\ref{thm:chain} at \(p\), with the baseline truncated at zero and the zero-gain and dropped primes treated as in \cite[Def.~12.3 and \S12.1]{Adr}. The global gain and the remaining index correction are
\[
 \Gamma=\sum_{p\in\mathcal P}2\log p\left(\sum_{j=1}^m\mathrm{gain}^{\mathcal D}_j(p)-\tau(p)\right),
 \qquad I=\IO=\sum_{p\in\mathcal P}\io(p)\log p .
\]
Here \(\mathcal P\) is the set of admitted inert primes of \cite[\S12.1]{Adr}. The uniform scaling \(\tau\) has already been deducted in \(\Gamma\).  As in Chapter~\ref{as-chap}, \(\Gamma\) is the total gain; it is not the Galois group of Chapter~\ref{gf-chap}.  The two theorems below use the following \emph{global inputs}, which \cite[\S12.6, (L1)--(L8)]{Adr} takes from the constructions of the lifts and of the auxiliary prime: at each served place, the local lattices, Hecke operators and orbit bounds of Theorem~\ref{thm:colength}; the base lattice \(M_u\) and its saturated ramified part; the auxiliary prime \(T_u\ge X'\), the nodes and their weights; a sublattice at \(T_u\) of colength at most \(\lambda_T\) times the rank, with \(\lambda_T\) as in Theorem~\ref{thm:K2}; and finite supports for the local terms, with residue characteristics whose logarithms sum to at most \(X''\).

\begin{theorem}[Prime sum (A1)]\label{thm:K1}
Assume the global inputs, the size conditions of \cite[\S12.6]{Adr}, and weights in the prescribed windows with \(\mathrm{lhsW}\le x^{-9}\). Here \(\mathrm{lhsW}\) sums \( (2\mathbf1_{\Sig(p)>x^{-\varepsilon}}+4\Sig(p))\log p/p\) over candidate primes. Then
\[
 S-\Gamma\le(2\log4+1)N_Tmb,\qquad 0\le I\le N_Tmb
\]
\cite[Thm.~12.27]{Adr}. The count of tops on orbits that are not good comes from \cite[Prop.~5.20]{Adr}; choosing weights with small \(\mathrm{lhsW}\) uses the engines.
\end{theorem}
\noindent \emph{Status}: derived in \cite[Thm.~12.27]{Adr} from the engines \cite[Hyp.~12.12]{Adr} and from the global inputs, which rest there on (N2) and (N3) through Theorem~\ref{thm:true-lifts}; these hypotheses are assumed there.  The formalization proves the final theorem by a different arrangement at this step (Chapter~\ref{gl-chap}).  Its counterpart of this theorem is Theorem~\ref{gl-thm:linkG}, an inequality between explicit numbers at each prime: the gains, prices and colengths of Theorem~\ref{thm:chain} and Theorem~\ref{thm:colength} against the prime sum.  It is proved with the chain of lifts replaced by the zero chain, for which the divisibility clause holds trivially (Problem~\ref{gl-pr:zero-chain}).  The divisibility of the lifts is Theorem~\ref{gl-thm:gain}, proved on a separate branch, where (a)--(c) of Theorem~\ref{thm:true-lifts} enter (Chapters~\ref{nk-chap} and~\ref{fo-chap}); the two branches meet only in the source lattice \(N_u\) (Section~\ref{gl-sec:record}).  The engines are not separate statements in the formalization.  Simplest incarnation: Theorem~\ref{ip-thm:main} supplies the positive parts for one inert prime; Theorem~\ref{od-thm:mod4} sums \(\log p/p\) over \(p\equiv3\bmod4\) with main term \(\frac12\log X\).

\begin{problem}\label{pr:offset}
The node offset costs \(b'/x\) per top and step at served primes \(p<X'=xb/b'\). From \(\sum_{p\le X}\log p\le X\log4\), show that its total cost is at most \(2\log4\,N_Tmb\). If the other losses together cost at most \(N_Tmb\), this gives the first inequality of Theorem~\ref{thm:K1}.
\end{problem}
\begin{problem}\label{pr:mean}
Put \(a=xb\). At a served prime, \(S_v\le2mN_Ta/p\) by the local gain estimate in the proof of \cite[Thm.~12.27]{Adr}. Show that \(I\le amN_T\,\mathrm{lhsW}\le N_Tmb\,x^{-8}\). Thus the mean-share input controls the excess colength even after summing over primes.
\end{problem}
\begin{problem}\label{pr:place-count}
Of \(2n\) places, \(N_2=\delta n\) carry weight \(b\). Average their total weight over all places, then sum over tops and steps to obtain \(N_Tmb\delta/2\). For a prime \(p\) with \(\chi_d(p)=-1\), show that \(1-\chi_d(p)^i\) gives weight two for odd \(i\) and zero for even \(i\), independently of \(N_2\). Thus changing \(N_2\) changes only the first of these two counts.
\end{problem}

For each chain index \(u\), \(M_u\) is the original hermitian lattice and \(N_u\) is the rescaled source obtained using \(N\!P_w\) and the auxiliary lattice at \(T_u\) \cite[Def.~12.20]{Adr}. Write \(\mu,\mu_{\max}\) for mean and maximal slopes. In the chain record \(L\) is the common real shift and \(E_u=\sum_jE_{u,j}\) the sum of the step corrections.
\begin{theorem}[Slope (A2)]\label{thm:K2}
Under the global inputs and the size conditions, with \(\Sig(p)\le1/2\) at every served prime,
\begin{align*}
 \mu(M_u)+\Gamma+L-E_u-e_{2,u}&\le\mu(N_u),\\
 \mu_{\max}(N_u)&\le\mu_{\max}(M_u)+\Gamma+L-E_u .
\end{align*}
Here \(e_{2,u}=4I+4\lambda_T\log T_u+R_X+1\),
\(4\lambda_T=C_T N_Tm(a/k+b')\), and \(R_X\le6X''+2\log T_u\), with \(k\) the base rank \cite[Thm.~12.21]{Adr}.
\end{theorem}
\noindent \emph{Status}: derived in \cite[Thm.~12.21]{Adr} from the same hypotheses as Theorem~\ref{thm:K1}, assumed there.  The formalization proves the final theorem; there the slope bound comes from the source lattice \(N_u\) of Section~\ref{gl-sec:record}, on which the divisibility of the lifts makes the pairing integral, and the colength it costs (Problem~\ref{gl-pr:record}).  Simplest incarnation: Theorem~\ref{ch-thm:chain} bounds the sum of finite gains by the difference of the two end slopes plus the archimedean step heights.

\begin{problem}\label{pr:aux}
If \(\log T_u\le Cmy\), then the \(a/k\) part of \(4\lambda_T\log T_u\) is at most \(CC_T\,yN_Ta\,m^2/k\). Under \(N_Tm(a/k+b')\ge1\), show that
\[
 e_{2,u}\le {\rm Lo}_u
 =C_* I+C_*N_Tm(a/k+b')\log T_u+C_*X''+1
\]
for \(C_*\ge\max(6,C_T+2)\).
\end{problem}
\begin{problem}\label{pr:record}
Take \(e_L=0\) and enlarge \(C_*\) to at least \(2\log4+1\). Theorem~\ref{thm:K1} then gives both inequalities of (K1). Show that Theorem~\ref{thm:K2} and Problem~\ref{pr:aux} give (K2), whose upper bound allows an extra \(C_*mB_0\ge0\). Conversely, replacing \({\rm Lo}_u\) in its lower bound by a common \(e_2\ge{\rm Lo}_u\) gives the lower bound (A2) used in assembly \cite[Lem.~18.13, Prop.~18.14]{Adr}.
\end{problem}

\paragraph{Deliverable.}
Theorems~\ref{thm:K1} and~\ref{thm:K2} supply (K1) and (K2) of \cite[Def.~18.9]{Adr} with \(I=\IO\) and \(e_L=0\).  In \cite{Adr} they are derived from (N1)--(N3) and the engines of \cite[Hyp.~12.12]{Adr}, of which (N2), (N3) and the engines are assumed there.  The formalization proves the final theorem.  At this step it proves the inequality of Theorem~\ref{thm:K1} between explicit numbers, with the zero chain in place of the lifts, and the divisibility of the lifts on a separate branch, through the kernel law, the unit scalars and the dual frames of Theorem~\ref{thm:true-lifts} (Chapters~\ref{nk-chap}, \ref{fo-chap} and~\ref{gl-chap}); the Hecke image is treated in Chapters~\ref{ea-chap} and~\ref{co-chap}, and the engines are not separate statements in it; see Section~\ref{as-sec:map} of Chapter~\ref{as-chap}.  Alongside (N2), (N3) and the three engines, \cite[\S18.7]{Adr} singles out the comparisons between levels and the test systems, the Weyl lattice model in \cite[Cor.~12.65]{Adr}, the polynomial bound for \(c_G\) in \cite[Prop.~12.67]{Adr}, and the colength at the auxiliary prime among the global inputs.

\section*{Hints}
\subsection*{\S\ref{sec:minors}}
\begin{description}
\item[\ref{pr:cap}] Use multilinearity in the columns.
\item[\ref{pr:digits}] Count digits at the two positions separately.
\end{description}
\subsection*{\S\ref{sec:chain}}
\begin{description}
\item[\ref{pr:minkowski}] Sum a geometric series.
\end{description}
\subsection*{\S\ref{sec:image}}
\begin{description}
\item[\ref{pr:vertex}] Separate \(i=0\) from \(i\ge1\).
\item[\ref{pr:share}] Use additivity through \(p^\tau M^{\rm ram}\).
\end{description}

\begin{thebibliography}{Adr}
\bibitem[Adr]{Adr}
J.~Adriazola, \emph{The Alp\"oge proof of the Riemann hypothesis}, manuscript, revised 24 November 2026, \url{https://github.com/jimjamzola/RH-Alpoge-Reconstruction}.
\bibitem[Lean]{Lean} The Lean community, \emph{riemann-hypothesis}, a Lean formalization of Alp\"oge's argument, \url{https://github.com/leanprover-community/riemann-hypothesis}.
\bibitem[LT]{LT} J.~Lubin and J.~Tate, Formal complex multiplication in local fields, Ann. of Math. (2) \textbf{81} (1965), 380--387.
\end{thebibliography}
\end{document}
